\begin{frame}{4. Formalisme \textit{Exact Potential Game}}
Pemilihan portofolio diformulasikan sebagai \textbf{\textit{Exact Potential Game}}, di mana utilitas murni individu (Pemain $i$) mendefinisikan keputusan mengkoleksi saham $x_i \in \{0, 1\}$:

$$ u_i(x_i, x_{-i}) = x_i \left( \tilde{\mu}_i - \frac{\gamma}{2}\Sigma_{ii} - \gamma \sum_{j \neq i} \tilde{\Sigma}_{ij} x_j \right) $$

\textbf{Inti Teori Permainan (\textit{Game Theory}):}
Terdapat interaksi strategis persaingan utilitas antar-saham:
\begin{itemize}
    \item Kenaikan risiko variansi individu ($\Sigma_{ii}$) menurunkan \textit{payoff}.
    \item Jika dua saham memiliki redundansi/korelasi tinggi ($\tilde{\Sigma}_{ij} > 0$), utilitas akan saling memberikan \textbf{penalti marginal} (hukuman) saat keduanya dipilih secara simultan ($x_i=1, x_j=1$).
\end{itemize}
\end{frame}
